\documentclass[../../main.tex]{subfiles}
\begin{document}
\section{光滑奇异同调}\label{section:ch18-smooth-singular-homology}

奇异同调与 de Rham 上同调之间的联系，将通过在奇异链上积分微分形式来建立。更确切地说，给定流形 $M$ 中的一个奇异 $p$-单形 $\sigma$ 和 $M$ 上的一个 $p$-形式 $\omega$，我们希望用 $\sigma$ 把 $\omega$ 拉回，并在 $\Delta^p$ 上积分所得形式。然而这个做法有一个直接的问题：形式只能由光滑映射拉回，而奇异单形一般只是连续的。本节中我们通过证明奇异同调同样可以用光滑单形来计算，来克服这个问题。

若 $M$ 是一个光滑流形，$M$ 中的一个\emph{光滑 $p$-单形}是一个映射 $\sigma:\Delta^p\to M$，它在如下意义下光滑：在每一点附近都有一个光滑延拓。$C_p(M)$ 中由光滑单形生成的子群记为 $C_p^\infty(M)$，称为 $p$ 次\emph{光滑链群}。由于光滑单形的边界是光滑链，我们可以把 $M$ 的 $p$ 次\emph{光滑奇异同调群}定义为商群
\begin{align*}
H_p^\infty(M)=\frac{\operatorname{Ker}(\partial:C_p^\infty(M)\to C_{p-1}^\infty(M))}{\operatorname{Im}(\partial:C_{p+1}^\infty(M)\to C_p^\infty(M))}.
\end{align*}
包含映射 $\eta:C_p^\infty(M)\hookrightarrow C_p(M)$ 与边界算子交换，因而在同调上诱导一个映射 $\eta_*:H_p^\infty(M)\to H_p(M)$。

\begin{theorem}[光滑奇异同调与奇异同调]\label{theorem:ch18-smooth-vs-singular}
对任意光滑流形 $M$，由包含诱导的映射 $\eta_*:H_p^\infty(M)\to H_p(M)$ 是一个同构。
\end{theorem}

证明的基本想法是：借助 Whitney 逼近定理，构造两个算子——一个光滑化算子 $s:C_p(M)\to C_p^\infty(M)$，以及一个同伦算子。细节相当技术性，我们在此只给出梗概。关键是对每个连续单形系统地构造一个到光滑单形的同伦，并使它与到 $\Delta^p$ 各边界面的限制相容。

\begin{figure}[ht]
\centering
\includegraphics[width=0.6\textwidth]{fig-18-4.png}
\caption{同伦 $H_\sigma$}\label{figure:ch18-homotopy-hsigma}
\end{figure}

\begin{lemma}\label{lemma:ch18-smooth-homotopy-construction}
设 $M$ 是一个光滑流形。对每个整数 $p\geqslant0$ 和每个奇异 $p$-单形 $\sigma:\Delta^p\to M$，存在一个连续映射 $H_\sigma:\Delta^p\times I\to M$，满足：
\begin{enumerate}[(i)]
\item $H_\sigma$ 是从 $\sigma(x)=H_\sigma(x,0)$ 到某个光滑 $p$-单形 $\widetilde\sigma(x)=H_\sigma(x,1)$ 的同伦；
\item 对每个面映射 $F_{i,p}$，$H_\sigma\circ F_{i,p}=H_\sigma\circ(F_{i,p}\times\operatorname{Id}_I)$；
\item 若 $\sigma$ 已是光滑 $p$-单形，则 $H_\sigma$ 是恒定同伦 $H_\sigma(x,t)=\sigma(x)$。
\end{enumerate}
\end{lemma}

\begin{figure}[ht]
\centering
\includegraphics[width=0.55\textwidth]{fig-18-5.png}
\caption{从 $\Delta^p\times I$ 到 $S$ 的收缩}\label{figure:ch18-retraction}
\end{figure}

\begin{figure}[ht]
\centering
\includegraphics[width=0.55\textwidth]{fig-18-6.png}
\caption{说明 $f$ 在 $\partial\Delta^p$ 上光滑}\label{figure:ch18-smooth-on-boundary}
\end{figure}

\begin{proof}
对 $\sigma$ 的维数作归纳构造同伦 $H_\sigma$。对 $0$-单形直接取 $H_\sigma(x,t)=\sigma(x)$。归纳步骤中，若 $\sigma$ 不光滑，则在“棱柱” $\Delta^p\times I$ 的底与侧面上先定义好映射，再用一个从 $\Delta^p\times I$ 到底面与侧面之并 $S$ 的收缩延拓到整个棱柱；最后利用 Whitney 逼近定理把顶端替换为一个光滑单形，并把两段同伦拼接起来。
\end{proof}

\begin{proof}[定理~\cref{theorem:ch18-smooth-vs-singular} 的证明]
用引理~\cref{lemma:ch18-smooth-homotopy-construction} 定义光滑化算子 $s\sigma=H_\sigma\circ i_1$。利用面映射的相容关系可验证 $s$ 是一个链映射。再构造一个同伦算子 $h$ 满足 $\partial h+h\partial=\eta\circ s-\operatorname{Id}$，由此推出 $\eta_*\circ s_*=\operatorname{Id}$ 且 $s_*\circ\eta_*=\operatorname{Id}$。因此 $\eta_*$ 是同构。
\end{proof}
\end{document}
